YS II.44

    svadhyayad ista-devata-samprayoga

YS II.45

    samadhi-siddhi îsvara-pranidhana

SK 22.

Prakriti mahat tata ahamkara tasmat ca sodasaka gana
sodasakat tasmad api pancabhyah panca bhutani

From nature evolves the Great Principle;
from this evolves the I Principle;
from this evolves the set of sixteen;
from the five of this set of sixteen evolves the five elements.

SK 23.

buddhi adhyavasaya dharma jnana viraga aisvarya
etad rupam sattvikam tasmat viparyastam tamasam

Buddhi is self-certainty.
Virtue, knowledge, dispassion, and power are its manifestations
when the sattva attribute abounds.
And the reverse of these, when the tamasa attribute abounds.

SK 24.

ahamkara abhimana tasmat pravartate dvividha sarga eva
ekadasaka gana ca tanmatra pancaka

Ahamkara is self-assertion;
from that proceeds a two-fold evolution
the set of eleven and the five-fold primary elements.

YS II.46

    sthira-sukham asanam

YS II.47

    prayatna-saithilya-ananta-samapattibhyam

YS II.48

    tato dvandva-anabhighata

SK 25.

vaikrita ahamkara ekadasaka sattvika pravartate
tanmatra bhutade sa tamasa taijasa ubhayam

The set of eleven sattvic indriyas proceeds from the Vaikriti form of I-Principle;
the set of five tamasic tanmatras proceed from the Bhutadi form of I-Principle.
From the Taijasa form of I-Principle proceed both of them.

SK 26.

buddhi indriyani akhyani caksu srotram ghrana rasana tvak
karma indriyani ahu vak pani pada payu upastha

Organs of knowledge are called the Eye, the Ear, the Nose, the Tongue, the Skin.
Organs of action are called the Speech, the Hand, the Feet, the excretory organ and the organ of generation.

YS II.49

    tasmin sati svasa-prasvasayo gati-viccheda pranayama

YS II.50

    bahya-abhyantara-stambha-vritti desa-kala-sankhyabhi-paridrsta dirgha-suksma

YS II.51

    bahya-abhyantara-visaya-aksepi caturtha

SK 27.

atra mana ubhayatmakam sankalpakam ca sadharmaya indriyam
nanatvam bahya bheda ca guna-parinama-visesa

Of these (sense organs), the Mind possesses the nature of both.
It is the deliberating principle, and is also called a sense organ
since it posseses properties common to the sense organs.
Its multifariousness and also its external diversities are owing
to special modifications of the Attributes.

SK 28.

pancanam rupadisu alocanamatram isyate
pancanam vritti vacana adana viharana utsarga ca ananda

The function of the five in respect to form and the rest,
is considered to be mere observation.
Speech, manipulation, locomotion, excretion and gratification
are the functions of the other five.

YS II.52

    tata ksiyate prakasa-avaranam

YS II.53

    dharanasu ca yogyata manasa

YS II.54

    sva-visaya-asamprayoge cittasya svarupa-anukara ive indriyanam pratyahara

YS II.55

    tata parama vasyata indriyanam

YS III.1

    desa-bandha cittasya dharana

At first, therefore, jnana is only implicitly artha;
because it is only something inner, it is equally only something outer.
It is at first simply an immediate and in this shape its moments have
the form of immediate, fixed determinations.
It appears as the determinate jnana,
as the sphere of mere Manas.
– Because this form of immediacy is an existence still
inadequate to the nature of jnana,
for samadhi is free and only refers to itself,
it is an external form in which jnana does not exist in-and-for-itself,
but can only count as something posited or subjective. –
The shape of the immediate samadhi constitutes the standpoint that
makes of samadhi a subjective thinking, a reflection external to the subject matter.
This stage constitutes, therefore, dharana, or the formal samadhi.
Its externality is manifested in the fixed being of its determinations that
makes them come up each by itself, isolated and qualitative,
and each only externally referred to the other.
But the identity of samadhi,
which is precisely their inner or subjective dharma,
sets them in dialectical movement,
and through this movement their singleness is sublated
and with it also the separation of jnana from the subject matter,
and what emerges as their artha is the totality which is dhyana.

YS III.2

tatra pratyaya-eka-tanata dhyanam

Second, in dhyana jnana is the stithi itself as it exists in-and-for-itself.
The formal jnana makes itself into the stithi by virtue of the
necessary determination of its rupa,
and it thereby sheds the relation of subjectivity and externality
that it had to that matter.
Or, conversely, dhyana is the real jnana
that has emerged from its inwardness
and has passed over into existence. –
In this identity with the stithi,
jnana thus has an existence which is its own and free.
But this existence is still a freedom which is immediate and not yet negative.
Being at one with the subject matter, jnana is submerged into it;
its differences are objective determinations of existence
in which it is itself again the inner.
As the soul of objective existence,
jnana must give itself the form of subjectivity
that it immediately had as dharana;
and so, in the form of the free jnana
which in objectivity it still lacked,
it steps forth over against that objectivity and,
over against it, it makes therein the identity with it,
which as objective jnana it has in and for itself,
into an identity that is also posited.


YS III.3

    tad evartha-matra-nirbhasa svarupa-sunyam iva samadhi

In this consummation in which jnana has
the form of freedom even in its dhyana,
the adequate jnana is samadhi.
Buddhi, which is the sphere of samadhi,
is the self-unveiled artha in which jnana attains
the realization absolutely adequate to it,
and is free inasmuch as in this real world,
in its dhyana, it recognizes its dharana,
and in this dharana recognizes that dhyana.

YS III.4

    trayam ekatra samyama

YS III.5

    taj-jayat prajna-aloka

YS III.6

    tasya bhumisu viniyoga

YS III.7

    trayam antar-angam purvebhya

YS III.8

    tad api bahir-angam nirbijasya

SK 29.

trayasya svalaksanyam vritti sa esa asamanya bhavati
samanya karana-vritti pranadhyah-vayavah-panca

Of the three internal instruments, their own characteristics are their functions;
this is peculiar to each.
The common modification of the instruments is the five airs
such as prana and the rest.

YS III.9

    vyutthana-nirodha-samskarayor abhibhava-pradur-bhavau nirodha-kshana-cittanvayo nirodha-parinama

YS III.10

    tasya prasanta-vahita samskarat

    Manas, the understanding, first the faculty for the cognition of the general (of rules)

    a. Concept as such

    The concept is, to start with, formal, the concept in its beginning or as
    the immediate concept. – In this immediate unity, its difference or its
    positedness is, first, itself initially simple and only a reflective shine, so that
    the moments of the difference are immediately the totality of the concept
    and only the concept as such.
    But, second, because it is absolute negativity, the concept divides and
    posits itself as the negative or the other of itself; yet, because it is still
    immediate concept, this positing or this differentiation is characterized by
    the reciprocal indifference of its moments, each of which comes to be on
    its own; in this division the unity of the concept is still only an external
    connection. Thus, as the connection of its moments posited as self-subsisting
    and indifferent, the concept is judgment. 29
    Third, although the judgment contains the unity of the concept that
    has been lost in its self-subsisting moments, this unity is not posited. It
    will become posited by virtue of the dialectical movement of the judgment
    which, through this movement, becomes syllogistic inference, 30 and this is
    the fully posited concept, for in the inference the moments of the concept
    as self-subsisting extremes and their mediating unity are both equally posited.
    But since this unity itself, as unifying middle, and the moments, as self-
    subsisting extremes, stand at first immediately opposite one another, this
    contradictory relation that occurs in the formal inference sublates itself, and
    the completeness of the concept passes over into the unity of totality; the
    subjectivity of the concept into its objectivity.

SK 30.

drste catustayasya tu yugapat vrttih tasya kramasasca nirdista
tatha api adrste trayasya vrttih tat purvika

Of all the four, the functions are said to be simultaneous and also successive
with regard to the seen objects;
with regard to the unseen objects
the functions of the three are preceeded by that.

The concept, as this connection of its self-subsistent determinations, has
lost itself, for the concept itself is no longer the posited unity of these deter-
minations, and these no longer are moments, the reflective shining of the
concept, but subsist rather in and for themselves. – As singularity, the con-
cept returns in determinateness into itself, and therewith the determinate
has itself become totality. The concept’s turning back into itself is thus the
absolute, originative partition of itself, that is, as singularity it is posited as
judgment.

    sarva-arthata-ekagratayo kshayodayau cittasya samadhi-parinama

    The logical form of all judgments consists of the objective unity of the apperception of the concepts contained therein

    Judgment

    Ahamkara, the determinative power of judgment,
    second the faculty for the subsumption of the particular under the general

Judgment is the determinateness of the concept posited in the concept itself.
The determinations of the concept, or, what amounts to the same thing
as shown, the determinate concepts, have already been considered on
their own; but this consideration was rather a subjective reflection or a
subjective abstraction. But the concept is itself this act of abstracting;
the positioning of its determinations over against each other is its own
determining. Judgment is this positing of the determinate concepts through
the concept itself.

Judging is therefore another function than conceiving; or rather, it is the
other function of the concept, for it is the determining of the concept through
itself. The further progress of judgment into a diversity of judgments is
this progressive determination of the concept. What kind of determinate
concepts there are, and how they prove to be necessary determinations of
it – this has to be exhibited in judgment.

Judgment can therefore be called the first realization of the concept,
for reality denotes in general the entry into existence as determinate being.
More precisely, the nature of this realization has presented itself in such
a way that the moments of the concept are totalities which, on the one
hand, subsist on their own through the concept’s immanent reflection or
through its singularity; on the other hand, however, the unity of the concept
is their connection. The immanently reflected determinations are determi-
nate totalities that exist just as essentially disconnected, indifferent to each
other, as mediated through each other. The determining itself is a totality
only as containing these totalities and their connections. This totality is the
judgment. – The latter contains, therefore, the two self-subsistents which
go under the name of subject and predicate. What each is cannot yet be said;
they are still indeterminate, for they are to be determined only through the
judgment. Inasmuch as judgment is the concept as determinate, the only
determination at hand is the difference that it contains between determinate
and still indeterminate concept. As contrasted to the predicate, the subject
can at first be taken, therefore, as the singular over against the universal, or
also as the particular over against the universal, or the singular over against
the particular; so far, they stand to each other only as the more determinate
and the more universal in general.

It is therefore fitting and unavoidable to have these names, “subject”
and “predicate,” for the determinations of the judgment; as names, they
are something indeterminate, still in need of determination, and therefore
nothing but names. It is partly for this reason that the determinations
themselves of the concept could not be used for the two sides of judgment;
but a still stronger reason is because of the nature of a concept determination
which is nothing abstract, nothing fixed, but contains its opposite in it,
explicitly posited there; since the sides of the judgment are themselves
concepts and therefore the totality of the determinations of the concept,
each side must run through all these determinations, exhibiting them
within whether in abstract or concrete form. But now, if in this altering of
determination we want to fix the two sides in some general way, names will
be the most useful means, for they can be kept the same throughout the
process. – But a name remains distinct from the fact or the concept, and this
is a distinction that transpires within the judgment as such; since the subject
is in general the determinate term and more, therefore, of an immediate
existent, whereas the predicate expresses the universal, the essence or the
concept, the subject as such is at first only a kind of name; what it is, is first
enunciated only by the predicate which contains being in the sense of the
concept. When we ask, “What is this?,” or “What kind of plant is this?,”
the being we are enquiring about is often just a name, and once we learn
this name, we are satisfied that we now know what the fact is. This is being
in the sense of the subject. The concept, however, or at least the essence and
the universal in general, is only given by the predicate, and when we ask for
it, we do it in the sense of the judgment. – God, therefore, or spirit, nature,
or what have you, is as the subject of a judgment only a name at first; what
any such subject is in accordance with the concept, is first found only in
the predicate. When we ask for the predicate that belongs to such subjects,
the required judgment must be based on a concept that is presupposed; yet
it is the predicate that first gives this concept. It is, therefore, the mere
representation that in fact makes up the presupposed meaning, and this
yields only a nominal definition whereby it is a mere accident, a historical
fact, what is understood by a name. So many disputes about whether a
predicate does or does not belong to a subject are, therefore, nothing more
than verbal disputes, for they proceed from this form; what lies at the base
is still nothing more than a name.

Secondly, we now have to examine more closely how the connection of
subject and predicate in judgment is determined, and how the two are
themselves thereby determined. Judgment has in general totalities for its
sides, totalities that are at first essentially self-subsistent. The unity of the
concept is at first, therefore, only a connection of self-subsistent terms; it is
not yet the concrete, the fulfilled unity that has returned into itself from this
reality but is a unity rather outside which the two terms persist as extremes
yet unsublated in it. – Now any consideration of the judgment can start
either from the originative unity of the concept or from the self-subsistence
of the extremes. Judgment is the self-diremption of the concept; therefore,
it is by starting from the unity of the concept as ground that the judgment is
considered in accordance with its true objectivity. In this respect, judgment
is the originative division (or Teilung, in German) of an originative unity; the
German word for judgment, Urteil (or “primordial division”), thus refers
to what judgment is in and for itself. But the concept is present in the
judgment as appearance, since its moments have attained self-subsistence
there, and it is to this side of externality that ordinary representation is more
likely to fasten.

From this subjective standpoint, the subject and the predicate are there-
fore treated as ready-made, each for itself outside the other – the subject as
a subject matter that would exist even if it did not have that predicate, and
the predicate as a universal determination that would exist even without
accruing to this subject. The act of judgment accordingly brings with it
the further reflection whether this or that predicate which is in someone’s
head can and should be attached to the subject matter that exists outside it
on its own; the judgment itself is simply the act that combines the predicate
with the subject, so that, if this combination did not occur, the subject
and predicate would still each remain what it is, the one concretely existing
as thing in itself, the other as a representation in someone’s head. – But
the predicate which is combined with the subject should also pertain to it,
which is to say, should be in and for itself identical with it. The significance
of their being combined is that the subjective sense of judgment, and the
indifferent external persistence of the subject and predicate, are again sub-
lated. Thus in “this action is good,” the copula indicates that the predicate
belongs to the being of the subject and is not merely externally combined
with it. Of course, grammatically speaking this kind of subjective relation
that proceeds from the indifferent externality of subject and predicate is
perfectly valid, for it is words that are here externally combined. – It can also
be mentioned in this context that a proposition can indeed have a subject
and predicate in a grammatical sense without however being a judgment
for that. The latter requires that the predicate behave with respect to the
subject in a relation of conceptual determination, hence as a universal with
respect to a particular or singular. And if what is said of a singular sub-
ject is itself only something singular, as for instance, “Aristotle died at the
age of 73 in the fourth year of the 115th Olympiad,” then this is a mere
proposition, not a judgment. There would be in it an element of judgment
only if one of the circumstances, say, the date of death or the age of the
philosopher, came into doubt even though the stated figures were asserted
on the strength of some ground or other. In that case, the figures would be
taken as something universal, as a time that, even without the determinate
content of Aristotle’s death, would still stand on its own filled with some
other content or simply empty. Likewise would the news that my friend
N. has died be a proposition, and a judgment only if there were a question
as to whether he is actually dead and not just apparently dead.

In the usual definition of judgment, that it is the combination of two
concepts, we may indeed accept the vague expression of “combination” for
the external copula, and also accept that the terms combined are at least
meant to be concepts. But the definition is otherwise a highly superficial
one. It is not just that in the disjunctive judgment, for instance, there are
more than two so-called concepts that are combined; more to the point is
rather that the definition is much better than the matter defined, for it is not
determinations of concepts, but determinations of representation that are in
fact meant; it was remarked in connection with the concept in general, and
with the concept as determinate, that what usually goes under this name
of concept does not deserve the name at all; 49 where should concepts then
come from in the case of judgment? – Above all this definition of judgment
ignores what is essential to it, namely the difference of its determinations;
still less does it take into account its relation to the concept.

As regards the further determination of the subject and predicate, we
have remarked above that it is in judgment that they must first receive
their determination. But since judgment is the posited determinateness
of the concept, this determinateness possesses the given differences imme-
diately and abstractly as singularity and universality. – But inasmuch as
judgment is in general the immediate existence or the otherness of the con-
cept that has not yet restored itself to the unity through which it exists as
concept, there also emerges the determinateness that is void of concept,
the opposition of being and reflection or the in-itself. But since the concept
constitutes the essential ground of judgment, these determinacies are at least
indifferent in the sense that, when one accrues to the subject and the other
to the predicate, the converse relation equally holds. The subject, being
the singular, appears at first as the existent or as the one that exists for itself
with the determinate determinateness of a singular on which judgment is
passed – as an actual object even when it is such in representation only –
as for instance in the case of bravery, right, agreement, etc. The predicate,
which is the universal, appears on the contrary as the reflection of this
judgment on that object, or rather as the object’s immanent reflection that
transcends the immediacy of the judgment and sublates its determinacies
as mere existents – appears, that is, as the object’s in-itselfness. – In this way,
the start is made from the singular as the first, the immediate, and through
the judgment this singular is raised to universality, just as, conversely, the
universal that exists only in itself descends in the singular into existence or
becomes a being that exists for itself.

This significance of the judgment is to be taken as its objective meaning
and at the same time as the true significance of the previous forms of
transition. The existent comes to be and becomes another, the finite passes
over into the infinite and in it passes away; the existent comes forth into
appearance out of its ground and to this ground it founders; the accidents
manifest the wealth of substance as well as its might; in being, there is
transition into an other; in essence, there is the reflective shining in an
other that manifests the necessity of a connection. This transition and this
reflective shining have now passed over into the originative division of the
concept in judgment, and this division, in bringing the singular back to the
in-itselfness of its universality, equally determines the universal as something
actual. These two are one and the same – the positing of singularity in its
immanent reflection and of the universal as determinate.
But equally pertaining to this objective meaning is that the said differ-
ences, as they re-occur in the determinateness of the concept, are at the
same time posited as only appearing, that is to say, that they are nothing
fixed but accrue rather just as much to one determination of the concept as
to the other. The subject is therefore equally to be taken as the in-itself, and
the predicate as determinate existence in contrast to it. The subject without
the predicate is what the thing without properties, the thing-in-itself, is in
the sphere of appearance, an empty indeterminate ground; it is then the
implicit concept that receives a difference and a determinateness only in
the predicate; the predicate thus constitutes the side of the determinate
existence of the subject. Through this determinate universality the subject
refers to the outside, is open to the influence of other things and thereby
confronts them actively. What is there comes forth from its in-itselfness into
the universal element of combination and relations, into negative references
and into the interplay of actuality which is a continuation of the singular
into other singulars and is, therefore, universality.

Yet the identity just indicated, the fact that the determination of the
subject accrues equally to the predicate and vice versa, is not just a matter
for our consideration; it is not only in itself but is also posited in the
judgment; for the judgment is the reference connecting the two; the copula
expresses that the subject is the predicate. The subject is the determinate
determinateness, and the predicate is this determinateness of the subject as
posited; the subject is determined only in its predicate, or is subject only in
it; in the predicate, it is turned back into itself and is therein the universal. –
Now in so far as the subject is the self-subsistent term, this identity has the
relation that the predicate does not possess a self-subsistence of its own but
has its subsistence only in the subject; it inheres in the subject. Accordingly,
since the predicate is distinguished from the subject, it is only a singularized
determinateness of the subject, only one of its properties; the subject itself
is however the concrete, the totality of manifold determinacies, just as the
predicate contains one of them; the subject is the universal. – But, on the
other hand, the predicate also is self-subsistent universality, and the subject
conversely only one determination of it. The predicate thus subsumes the
subject; the singularity and the particularity are not for themselves but
have their essence and their substance in the universal. The predicate
expresses the subject in its concept; the singular and the particular are to the
subject accidental determinations; the subject is their absolute possibility.
When by “subsumption” an external connection of subject and predicate
is thought, and the subject is represented as something self-subsistent,
then subsumption refers to the subjective act of judging mentioned above,
namely the judging that starts off from the self-subsistence of both subject
and predicate. 51 Subsumption is then only the application of the universal to
a particular or singular posited under it in accordance with an indeterminate
representation, one of lesser quantity.

When we treat the identity of subject and predicate as meaning that at
one time one determination of the concept belongs to the subject and the
other to the predicate, and at another time the converse equally applies,
then the identity is as yet still implicit; on account of the self-subsistent
diversity of the two sides of judgment, their posited connection also has
the two at first as diverse. But it is the identity void of difference that in fact
constitutes the true connection of the subject and predicate. The determi-
nation of the concept is itself essentially a connection, for it is a universal; the
same determinations, therefore, which the subject and the predicate each
have, are also had by their connection. The connection is universal, for it
is the positive identity of both, of the subject and predicate; but it is also
determinate, for the determinateness of the predicate is the determinateness
of the subject; it is singular as well, for in it the self-subsisting extremes are
sublated as in their negative unity. – In judgment, however, this identity
is not posited yet; the copula is as the still indeterminate connection of
being in general, “A is B,” for the self-subsistence of the concept’s deter-
minacies, or the extremes, is in judgment the reality that the concept has
within. If the “is” of the copula were already posited as the determinate
and fulfilled unity of subject and predicate earlier mentioned, 52 were
posited as their concept, it would then already be the conclusion of syllogistic
inference.

To restore again this identity of the concept, or rather to posit it – this
is the goal of the movement of the judgment. What is already present in
the judgment is, on the one hand, the self-subsistence but also reciprocal
determinateness of the subject and predicate, and, on the other hand,
their still abstract connection. “The subject is the predicate” – this is what
the judgment says at first. But since the predicate is not supposed to be
what the subject is, a contradiction is at hand that must resolve itself, must
pass over into a result. Or rather, since the subject and predicate are in
and for themselves the totality of the concept, and judgment is the reality
of the concept, the judgment’s forward movement is only development;
what comes forth from it is already present in it, and to this extent the
demonstration is a display, 54 a reflection as the positing of that which is already
at hand in the extreme terms of the judgment; but even this positing is
already present; it is the connection of the extremes.

SK 31.

svam svam vrittim pratipadyante paraspara akuta hetukam
purushartha eva hetu na kenacit karanam karyate

Instruments enter into their respective functions being incited by mutual impulse.
The meaning of the Spirit is the sole motive (for the activities of the instruments).
By none whatsoever is an instrument made to act.

On closer examination, the positive factor in this result
which is responsible for the transition of the judgment into
another form is that, as we have just seen,
the subject and predicate are in the apodictic judgment each the whole concept.
The unity of the concept, as the determinateness constituting
the copula that connects them, is at the same time distinct from them.
At first, it stands only on the other side of the subject as
the latter’s immediate constitution.
But since its essence is to connect,
it is not only that immediate constitution but
the universal that runs through the subject and predicate.
While subject and predicate have the same content,
it is the form of their connection that is instead posited
through the determinateness of the copula,
the determinateness as a universal or the particularity.
Thus it contains in itself both the form determinations of
the extremes and is the determinate connection of the subject and predicate:
the accomplished copula of the judgment, the copula replete of content,
the unity of the concept that re-emerges from the judgment
wherein it was lost in the extremes.
By virtue of this repletion of the copula, the judgment has become syllogism.

YS III.12

    tata punashantoditau tulya-pratyayau cittasya-ekagrata-parinama

    Syllogism

    Buddhi, reason, third the faculty for the determination of
    the particular from the general (for the derivation from principles)

The syllogism is the result of the restoration of the concept in the judgment,
and consequently the unity and the truth of the two. The concept as such
holds its moments sublated in this unity; in judgment, the unity is an
internal or, what amounts to the same, an external one, and although the
moments are connected, they are posited as self-subsisting extremes. In the
syllogism, the determinations of the concept are like the extremes of
the judgment, and at the same time their determinate unity is posited.
Thus the syllogism is the completely posited concept; it is, therefore, the
rational. – The understanding is taken to be the faculty of the determinate
concept which is held fixed for itself by virtue of abstraction and the form
of universality. But in reason the determinate concepts are posited in their
totality and unity. Therefore, it is not just that the syllogism is rational but
that everything rational is a syllogism. Syllogistic inference has long since
been ascribed to reason; but, on the other hand, reason in and for itself,
and rational principles and laws, are so spoken of that no light is thrown
on why the one reason that syllogizes, and the other which is the source
of laws and otherwise eternal truths and absolute thoughts, hang together.
If the former is supposed to be only a formal reason while the latter is
supposed to be the one that generates content, then one would expect on
this distinction that precisely the form of reason, the inference, would not
be missing in the latter. And yet, the two are commonly held so far apart,
the one without mention of the other, that it seems as though the reason
of absolute thoughts were ashamed, so to speak, of inferential reason, and
the syllogism were listed as also an activity of reason merely as matter of
tradition. But surely, as we have just remarked, logical reason must be
essentially recognizable, when regarded as formal, also in the reason that
is concerned with a content; indeed, no content can be rational except
by virtue of the rational form. In this matter we cannot rely on what is
commonly said about reason, for common views fail to tell us what we are
to understand by reason; this would-be rational wisdom is so busy with its
objects that it forgets to pay attention to reason itself but only identifies it
by characterizing it through the objects that it is said to have. If reason is
supposed to be a cognition that would know about God, freedom, right and
duty, the infinite, the unconditional, the suprasensible, or even gives only
representations and feelings of such objects, then for one thing these objects
are only negative, and for another the original question still stands, what
is there in all these objects that makes them rational? – The answer is that
the infinitude in them is not the empty abstraction from the finite, is not a
universality which is void of content and determination, but is the fulfilled
universality, the concept which is determined and is truly in possession of its
determinateness, namely, in that it differentiates itself internally and is the
unity of its thus intelligible and determined differences. Only in this way
does reason rise above the finite, the conditioned, the sensuous, or however
one might define it, and is in this negativity essentially replete with content,
for as unity it is the unity of determinate extremes. And so the rational is
nothing but the syllogism.

Now the syllogism, like judgment, is at first immediate; as such, its
determinations (termini) are simple, abstract determinacies; it is then the
syllogism of the understanding. If one stays at this configuration of the
syllogism, then its rationality, though present there and posited, is not
apparent. The essential element of the syllogism is the unity of the extremes,
the middle term that unites them and the ground that supports them.
Abstraction, by holding fast to the self-subsistence of the extremes, posits this
unity opposite them, as a determinateness with just as fixed an existence of its
own, thus grasping it more as a non-unity than as a unity. The expression,
“middle term” (medius terminus), is derived from spatial representation,
and has its share of responsibility for why we stop short at the externality
of the terms. Now if the syllogism consists in the positing in it of the unity
of the extremes, but if this unity is simply taken on the one hand as a
particular by itself, and on the other hand as only an external connection,
and non-unity is made the essential relation of syllogism, then the reason
of the syllogism is of no help to rationality.

SK 32.

karanam trayodasavidham tad aharana dharana prakasakaram
tasya karyam ca dasadha aharyam dharyam prakasyam ca

Instruments are of thirteen kinds performing the functions of
seizing, sustaining and illuminating.
Its objects are of ten kinds, vis
the seized, the sustained, and the illumined.

First, the syllogism of existence, in which the terms are thus
immediately and abstractly determined, demonstrates internally that,
since like judgment it is the connection of those terms,
these are not in fact abstract but each contains in it
the reference connecting it to the others,
and the determination of the middle term is not just
a determinateness opposed to the determinations of the extremes
but contains these extremes posited in it.

Through this dialectic, the syllogism of existence becomes
the syllogism of reflection,  the second syllogism.
Its terms are such that in each the other shines
essentially reflected in it, or are posited as mediated,
as they are indeed supposed to be in accordance with
the nature of syllogistic inference in general.

Third, inasmuch as this reflective shining or this mediatedness is reflected
into itself, syllogism is determined as the syllogism of necessity,
one in which the mediating factor is the objective nature of the fact.
As this syllogism determines the extremes of the concept also as totalities,
it has attained the correspondence of its concept (or the middle term) and
its existence (or the difference of the extremes).
It has attained its truth – and with that it has stepped forth
out of subjectivity into objectivity.

YS III.13

    etena bhutendriyesu dharma-laksana-vastha-parinama vyakhyata

    a. Mechanism

Since objectivity is the totality of the concept that has returned into its
unity, an immediate is thereby posited which is in and for itself that
totality, and is also posited as such, but in it the negativity of the concept
has as yet not detached itself from the immediacy of the totality; in other
words, the objectivity is not yet posited as judgment. In so far as it has
the concept immanent in it, the difference of the concept is present in it;
but on account of the objective totality, the differentiated moments are
complete and self-subsistent objects that, consequently, even in connection
relate to one another as each standing on its own, each maintaining itself
in every combination as external. – This is what constitutes the character
of mechanism, namely, that whatever the connection that obtains between
the things combined, the connection remains one that is alien to them,
that does not affect their nature, and even when a reflective semblance
of unity is associated with it, the connection remains nothing more than
composition, mixture, aggregate, etc. Spiritual mechanism, like its material
counterpart, also consists in the things connected in the spirit remaining
external to one another and to spirit. A mechanical mode of representation,
a mechanical memory, a habit, a mechanical mode of acting, mean that the
pervasive presence that is proper to spirit is lacking in what spirit grasps
or does. Although its theoretical or practical mechanism cannot take place
without its spontaneous activity, without an impulse and consciousness,
the freedom of individuality is still lacking in it, and since this freedom
does not appear in it, the mechanical act appears as a merely external
one.

a. the mechanical object

The object is, as we have seen, the syllogism whose mediation has attained
equilibrium and has therefore come to be immediate identity. It is therefore
in and for itself a universal – universality, not in the sense of a commonalit
of properties, but a universality that pervades particularity and in it is
immediate singularity.

b. the mechanical process

If objects are regarded only as self-enclosed totalities, they cannot act on
one another. Regarded in this way, they are the same as the monads which,
precisely for that reason, were thought of as having no influence on each
other. But the concept of a monad is for just this reason a deficient reflec-
tion. For, in the first place, the monad is a determinate representation of its
only implicit totality; as a certain degree of development and positedness of its
representation of the world, it is determinate; but since it is a self-enclosed
totality, it is also indifferent to this determinateness and is, therefore, not its
own determinateness but a determinateness posited through another object.
In second place, it is an immediate in general, for it is supposed to be just a
mirroring; 11 its self-reference is therefore abstract universality and hence an
existence open to others. – It does not suffice, in order to gain the freedom of
substance, to represent the latter as a totality that, complete in itself, would
have nothing to receive from the outside. On the contrary, a self-reference
that grasps nothing conceptually but is only a mirroring 11 is precisely a
passivity towards the other. – Likewise the determinateness, whether we now
take it as the determinateness of a being that exists or that mirrors, 11 as a
degree of the monad’s own internally generated development, is something
external; the degree that the development achieves has its limit in an other.
To project the reciprocal influence of substances into a predetermined har-
mony means nothing more than to make it a presupposition, in effect to
remove it from the scope of the concept. – The need to avoid the interac-
tion of substances was founded on the moment of absolute self-subsistence
and originariness which was made a fundamental assumption. But since
the positedness, the degree of development, does not correspond to this
assumed in-itselfness, it has for this reason its ground in an other.
In connection with the relation of substantiality, we showed that it passes
over into the relation of causality. 12 But now the existent no longer has the
determination of a substance but that of an object; the causal relation has
come to an end in the concept; the originariness of one substance vis-à-vis
another has shown itself to be a reflective shine, the substance’s action a
transition into the opposite substance. This relation has therefore no objec-
tivity. Hence in so far as one object is posited in the form of subjective
unity, as efficient cause, this no longer counts as an originary determination
but as something mediated; the active object has this determination only
by means of another object. – Mechanism, since it belongs to the sphere of
the concept, has that posited within it which proved to be the truth of the
relation of causality, namely, that the cause which is supposed to be some-
thing existing in and for itself is in fact effect just as well, positedness. In
mechanism, therefore, the originary causality of the object is immediately a
non-originariness; the object is indifferent to this determination attributed
to it; that it is a cause is therefore something accidental to it. – To this
extent, it can be said that the causality of substances is only the product
of representation. But precisely this causality as product of representation
is what mechanism is; for mechanism is this, that causality, as identical
determinateness of a diversity of substances and hence as the foundering
into this identity of their self-subsistence, is mere positedness; the objects
are indifferent to this unity and maintain themselves in the face of it. But
this also, their indifferent self-subsistence, is a mere positedness, and for this
reason they are capable of mixing and aggregating, and as an aggregate of
becoming one object. Through this indifference both to their transition and
to their self-subsistence, the substances are objects.

c. The product of the mechanical process

The product of formal mechanism is the object in general, an indifferent
totality in which determinateness is as posited. The object has hereby entered
the process as a determinate thing, and, in the extinction of this process,
the result is, on the one hand, rest, the original formalism of the object, the
negativity of its determinateness-for-itself. But, on the other hand, it is the
sublation of the determinateness, the positive reflection of it into itself, the
determinateness that has withdrawn into itself, or the posited totality of the
concept, the true singularity of the object. The object, determined at first in
its indeterminate universality, then as particular, is now determined as an
objective singular, so that in it that reflective semblance of singularity, which
is only a self-subsistence opposing itself to the substantial universality, is
sublated.
This resulting immanent reflection, the objective oneness of the objects,
is now a oneness which is an individual self-subsistence – the center.
Secondly, the reflection of negativity is the universality which is not a
fate standing over against determinateness, but a rational fate, immanently
determined – a universality that particularizes itself from within, the differ-
ence that remains at rest and fixed in the unstable particularity of the objects
and their process; it is the law. This result is the truth, and consequently
also the foundation, of the mechanical process.
c. absolute mechanism

SK 33.

antahkaranam trividham bahyam dasadha trayasya visayakhyam
bahyam sampratakalam abhyantaram karanam trikalam

The internal instruments are three-fold.
The external are ten-fold; they are called the objects of the three (internal instruments).
The external instruments function at the present time and
the internal instruments function at all the three times.

c. Transition of mechanism

This soul is however still immersed in its body. The now determined but
inner concept of objective totality is free necessity in the sense that the
law has not yet stepped in opposite its object; it is concrete centrality as
a universality immediately diffused in its objectivity. Such an ideality does
not have, therefore, the objects themselves for its determinate difference;
these are self-subsistent individuals of the totality, or also, if we look back
at the formal stage, non-individual, external objects. The law is indeed
immanent in them and it does constitute their nature and power; but
its difference is shut up in its ideality and the objects are not themselves
differentiated in the idealized non-indifference of the law. But the object
possesses its essential self-subsistence solely in the idealized centrality and its
laws; it has no power, therefore, to put up resistance to the judgment of the
concept and to maintain itself in abstract, indeterminate self-subsistence
and remoteness. Because of the idealized difference which is immanent in
it, its existence is a determinateness posited by the concept. Its lack of self-
subsistence is thus no longer just a striving towards a middle point, with
respect to which, precisely because its connection with it is only that of a
striving, it still has the appearance of a self-subsistent external object; it is
rather a striving towards the object determinedly opposed to it; and likewise
the center has itself for that reason fallen apart and its negativity has passed
over into objectified opposition. Centrality, therefore, is now the reciprocally
negative and tense connection of these objectivities. Thus free mechanism
determines itself to chemism.

YS III.14

    shantoditavyapadesya-dharmanupati dharmi

    Chemism

SK 34.

Tesam panca buddhi visesa-avisesa-visaya
Vak sabda-visaya-bhavati sesani tu pancavisayani

Of these, the five organs of knowledge have as their objects
both the specific and the general.
Speech has speech as its object;
the rest have all the five as its object.

YS III.15

    kramanyatvam parinamanyatve hetu

    Teleology

SK 35.

yasmat buddhi santahkarana sarvam visayam avagahate
tasmat trividham karanam dvari sesani dvarani

Since buddhi along with the other internal instruments
comprehends all of the objects these three instruments are
like the warders while the rest are like the gates.

We have now seen subjectivity, the being-for-itself of the concept, pass over
into the concept’s being-in-itself, into objectivity, and then the negativity of
that being-for-itself reassert itself in objectivity; the concept has so deter-
mined itself in that negativity that its particularity is an external objectivity,
or has determined itself as the simple concrete unity whose externality is its
self-determination. The movement of purpose has now attained this much,
namely that the moment of externality is not just posited in the concept,
the purpose is not just an ought and a striving, but as a concrete totality is
identical with immediate objectivity. This identity is on the one hand the
simple concept, and the equally immediate objectivity, but, on the other
hand, it is just as essentially mediation, and it is that simple immediacy only
through this mediation sublating itself as mediation. Thus the concept is
essentially this: to be distinguished, as an identity existing for itself, from
its implicitly existent objectivity, 44 and thereby to obtain externality, but in
this external totality to be the totality’s self-determining identity. So the
concept is now the idea.

YS III.16

    parinama-traya-samyamad atitanagata-jnanam

The idea of life has to do with a subject matter so concrete, and if you will so
real, that in dealing with it one may seem according to the common notion
of logic to have overstepped its boundaries. Needless to say, if the logic were
to contain nothing but empty, dead forms of thought, then there could be
no talk in it at all of such a content as the idea, or life, are. But if the subject
matter of logic is the absolute truth, and truth as such lies essentially in
cognition, then cognition at least would have to come in for consideration. –
It is common practice to have the so-called pure logic be followed by an
applied logic – a logic that has to do with concrete cognition, quite apart from
all the psychology and anthropology that is commonly deemed necessary to
interpolate into logic. But the anthropological and psychological side of
cognition is concerned with the form in which cognition appears when the
concept does not as yet have an objectivity equal to it, that is, when it
does not have itself as object. The part of the logic that deals with this
concrete cognition does not belong to applied logic as such; if it did, then
every science would have to be dragged into logic, for each is an applied
logic in so far as it consists in apprehending its subject matter in forms of
thought and of concepts. – The subjective concept has presuppositions that
are exhibited in psychological, anthropological, and other forms. But the
presuppositions of the pure concept belong in logic only to the extent that
they have the form of pure thoughts, of abstract essentialities, such as are the
determinations of being and essence. The same goes for cognition, which is
the concept’s comprehension of itself: no other shape of its presupposition
but the one which is itself idea is to be dealt with in the logic; this, however,
is a presupposition which is necessarily treated in logic. This presupposition
is now the immediate idea; for while cognition is the concept, in so far as
the latter exists for itself but as a subjectivity referring to an objectivity,
then the concept refers to the idea as presupposed or as immediate. But the
immediate idea is life.

To this extent the necessity of considering the idea of life in logic would
be based on the necessity, itself recognized in other ways, of treating the
concrete concept. But this idea has arisen through the concept’s own
necessity; the idea, that which is true in and for itself, is essentially the
subject matter of the logic; since it is first to be considered in its immediacy,
so that this treatment be not an empty affair devoid of determination, it is
to be apprehended and cognized in this determinateness in which it is life.
A comment may be in order here to differentiate the logical view of life
from any other scientific view of it, though this is not the place to concern
ourselves with how life is treated in non-philosophical sciences but only
with how to differentiate logical life as idea from natural life as treated in
the philosophy of nature, and from life in so far as it is bound to spirit. – As
treated in the philosophy of nature, as the life of nature and to that extent
exposed to the externality of existence, life is conditioned by inorganic nature
and its moments as idea are a manifold of actual shapes. Life in the idea is
without such presuppositions, which are in shapes of actuality; its presuppo-
sition is the concept as we have considered it, on the one hand as subjective,
and on the other hand as objective. In nature life appears as the highest stage
that nature’s externality can attain by withdrawing into itself and sublating
itself in subjectivity. It is in logic the simple in-itselfness which in the idea
of life has attained the externality truly corresponding to it; the concept
that came on the scene earlier as a subjective concept is the soul of life itself;
it is the impulse that gives itself reality through a process of objectification.
Nature, as it reaches this idea starting from its externality, transcends itself;
its end is not its beginning but is for it as a limit in which it sublates itself. –
Similarly, in the idea of life the moments of life’s reality do not receive the
shape of external actuality but remain enveloped in conceptual form.
In spirit, however, life appears both as opposed to it and as posited as at
one with it, in a unity reborn as the pure product of spirit. For life is here
to be taken generally in its proper sense as natural life, for what is called
the life of spirit as spirit, is spirit’s own peculiar nature that stands opposed
to mere life; just as we speak of the nature of spirit, even though spirit is
nothing natural but stands rather in opposition to nature. Thus life as such
is for spirit in one respect a means, and then spirit holds it over against
itself; in another respect, spirit is an individual, and then life is its body; in
yet another respect, this unity of spirit and its living corporeality is born
of spirit into ideality. None of these connections of life to spirit concerns
logical life, and life is to be considered here neither as the instrument of
a spirit, nor as a living body, nor again as a moment of the ideal and of
beauty. – In both cases, as natural life and as referring to spirit, life obtains a
determinateness from its externality, in one case through its presuppositions,
such as are other formations of nature, and in the other case through the
purposes and the activity of spirit. The idea of life by itself is free from both
the conditioning objectivity presupposed in the first case and the reference
to subjectivity of the second case.
Life, considered now more closely in its idea, is in and for itself absolute
universality; the objectivity which it possesses is throughout permeated by
the concept, and this concept alone it has as substance. Whatever is dis-
tinguished as part, or by some otherwise external reflection, has the whole
concept within it; the concept is the soul omnipresent in it, a soul which is
simple self-reference and remains one in the manifoldness that accrues to
the objective being. This manifoldness, as self-external objectivity, has an
indifferent subsistence which in space and time, if these could already be
mentioned here, is a mutual externality of entirely diverse and atomistic
matters. But externality is in life at the same time as the simple determi-
nateness of its concept; thus the soul flows omnipresently in this manifold
but remains at the same time the simple oneness of the concrete concept
with itself. – That way of thinking that clings to the determinations of
reflective relations and of the formal concept, when it comes to consider
life, the unity of its concept in the externality of objectivity, the absolute
multiplicity of atomistic matter, finds that all its thoughts are absolutely of
no avail; the omnipresence of the simple in the manifold externality is for
reflection an absolute contradiction and also, since it cannot at the same
time avoid witnessing this omnipresence in the perception of life and must
therefore grant the actuality of this idea, an incomprehensible mystery – for
reflection does not grasp the concept, nor does it grasp it as the substance
of life. – But this simple life is not only omnipresent; it is the one and
only subsistence and immanent substance of its objectivity; but as subjective
substance it is impulse, more precisely the specific impulse of particular dif-
ference, and no less essentially the one and universal impulse of the specific
that leads its particularization back to unity and holds it there. Only as this
negative unity of its objectivity and particularization is life self-referring, life
that exists for itself, a soul. As such, it is essentially a singular that refers to
objectivity as to an other, an inanimate nature. The originative judgment of
life consists therefore in this, that it separates itself off as individual subject
from the objective and, since it constitutes itself as the negative unity of
the concept, makes the presupposition of an immediate objectivity.
First, life is therefore to be considered as a living individual that is for itself
the subjective totality and is presupposed as indifferent to an objectivity
that stands indifferent over against it.
Second, it is the life-process of sublating its presupposition, of positing
as negative the objectivity indifferent to it, and of actualizing itself as the
power and negative unity of this objectivity. By so doing, it makes itself
into the universal which is the unity of itself and its other.
Third, consequently life is the genus-process, the process of sublating
its singularization and relating itself to its objective existence as to itself.
Accordingly, this process is on the one hand the turning back to its concept
and the repetition of the first forcible separation, the coming to be of a
new individuality and the death of the immediate first; but, on the other
hand, the withdrawing into itself of the concept of life is the becoming of
the concept that relates itself to itself, of the concept that exists for itself,
universal and free, the transition into cognition.

SK 36.

ete manas ahamkara pradipa-kalpa paraspara vilaksana
guna-visesa krstnam prakasya purushasya-artham buddhau prayacchati

These (external instruments along with manas and ahamkara)
which are characteristic-wise different from one another
and are different modifications of the attributes
and which resemble a lamp (in action) illuminating all (their respective objects)
present them to the Buddhi for the purpose of the Spirit,

The idea of spirit which is the subject matter of logic already stands,
on the contrary, inside pure science; it has no need, therefore, to observe
spirit’s tracing that course, to see how it gets entangled with nature, with
immediate determinateness, with matter, or in other words with pictorial
representation; this is what the other three sciences investigate. The idea of
spirit has this course already behind it, or what is the same, it has it rather
ahead of it – behind in so far as logic is taken as the final science; ahead
in so far as it is taken as the first science from which the idea first passes
over into nature. In the logical idea of spirit, therefore, the “I” is from the
start in the way it has emerged from the concept of nature as the truth of
nature, the free concept which in its judgment is itself the subject matter
confronting it, the concept as its idea. Also in this shape, however, the idea
is still not consummated.

Although the idea is indeed the free concept that has itself as its subject
matter, it is nonetheless immediate, and just because it is immediate, it is
still the idea in its subjectivity, and hence in its finitude in general. It is
the purpose that ought to realize itself, or the absolute idea itself still in its
appearance. What the idea seeks is the truth, this identity of the concept
itself and reality; but at first it only seeks it; for it is here as it is at first, still
something subjective. Consequently, although the subject matter that is for
the concept is here also a given subject matter, it does not enter into the
subject as affecting it, or as confronting it with a constitution of its own as
subject matter, or as a pictorial representation; on the contrary, the subject
transforms it into a conceptual determination; it is the concept which is the
active principle in it – which therein refers itself to itself, and, by thus
giving itself its reality in the object, finds truth.
Initially, therefore, the idea is one extreme of a syllogism, the concept
that as purpose has itself at first for its subjective reality; the other extreme is
the restriction of the subjective, the objective world. The two extremes are
identical in that they are the idea. Their unity is, first, that of the concept,
a unity which in the one extreme is only for itself and in the other only in
itself. Second, it is reality, abstract in the one extreme and in the other in
its concrete externality. – This unity is now posited through cognition, and,
because the latter is the subjective idea which as purpose proceeds from
itself, it is at first only a middle term. – The knowing subject, through the
determinateness of its concept which is the abstract being-for-itself, refers
to an external world; nevertheless, it does this in the absolute certainty
of itself, in order to elevate its implicit reality, this formal truth, to real
truth. It has the entire essentiality of the objective world in its concept; its
process consists in positing for itself the concrete reality of that world as
identical with the concept, and conversely in positing the latter as identical
with objectivity.The idea of cognition

Immediately, the idea of appearance is the theoretical idea, cognition as
such. For to the concept that exists for itself, the objective world immedi-
ately has the form of immediacy or of being, just as that concept is to itself
at first only the abstract concept of itself, is still shut up within itself. The
concept is therefore only as form, of which only its simple determinations
of universality and particularity are the reality that it possesses within, while
the singularity or the determinate determinateness, the content, is received
by it from the outside.

YS III.17

    sabda-artha-pratyayanam itaretara-adhyasat sankaras tat-pravibhaga-samyamat sarva-bhuta-ruta-jnanam

YS III.18

    samskara-saksat-karanat purva-jati-jnanam

YS III.19

    pratyayasya para-citta-jnanam

YS III.20

    na ca tat salambanam tasya-avisayi-bhutatvat

A. THE IDEA OF THE TRUE

At first the subjective idea is impulse. For it is the contradiction of the
concept that it has itself for the subject matter and is to itself the reality
without, however, the subject matter being an other that subsists on its
own over against it, or without the differentiation of itself from itself
having at the same time the essential determination of diversity and of
indifferent existence. The specific nature of this impulse is therefore to
sublate its own subjectivity, to make that first abstract reality a concrete
one, filling it with the content of the world presupposed by its subjectivity. –
From the other side, the impulse is determined in this way: the concept
is indeed the absolute certainty of itself; however, opposite its being-for-
itself there stands its presupposition of a world that exists in itself, but one
whose indifferent otherness has for the concept’s certainty of itself the status
of something merely unessential; the concept is therefore the impulse to
sublate this otherness and, in the object, to intuit its identity with itself.
This immanent reflection is the sublated opposition and the singularity
that originally makes its appearance as the presupposed being-in-itself but
is now posited and made actual for the subject; accordingly, this being-in-
itself is the self-identity of the form as it has issued from the opposition –
an identity which is therefore determined as indifferent towards the form
in its differentiation. It is content.
Consequently this impulse is the impulse of truth in so far as the truth
is in cognition, and therefore of truth in its strict sense as theoretical idea. –
Although objective truth is the idea itself as the reality that corresponds to
the concept, and to this extent a subject matter may or may not possess
truth, nevertheless the more precise meaning of truth is that it is such for
or in the subjective concept, in knowledge. 32 Truth is the relation of the
judgment of the concept, the concept that proved to be the formal judgment
of truth; 33 for the predicate in this judgment is not only the objectivity
of the concept, but the comparison connecting the concept of the fact with
the actuality of it. – This realization of the concept is theoretical in so far as
the concept still has, as form, the determination of subjectivity, or in so far
as it has for the subject the determination of being its own determination.
Because cognition is the idea as purpose or as subjective idea, the negation
of the world, presupposed as existing in itself, is the first negation; the
conclusion in which the objective is posited in the subjective has at first,
therefore, only the meaning that what exists in itself is posited only as
something subjective, only in conceptual determination, and consequently
does not exist as so posited in and for itself. Thus the conclusion only
attains to a neutral unity, or a synthesis, that is, to a unity of terms that
are originally separate, only externally conjoined. – Hence, since in this
cognition the concept posits the object as its own, the idea gives itself at first
only a content of which the foundation is given, in which only the form
of externality has been sublated. To this extent, this cognition still retains
its finitude in its realized purpose; in the realized purpose, it has at the
same time not attained its purpose, and in its truth it has not arrived at the
truth. For in so far as in the result the content still has the determination
of a given, the presupposed being-in-itself confronting the concept is not
sublated; the unity of concept and reality, the truth, is thereby equally not
contained in it. – Remarkable is that this side of finitude is the one that of
late has been clung to and accepted as the absolute relation of cognition –
as if the finite as such were to be the absolute! On this view, an unknown
thinghood-in-itself is attributed to the object, behind cognition, and this
thinghood, and the truth also along with it, are regarded for cognition as
an absolute beyond. Thought determinations in general, the categories, the
determinations of reflection, as well as the formal concept and its moments,
acquire on this view the status of determinations that are finite, not in and
for themselves, but in the sense of being something subjective as against
this empty thinghood-in-itself; the fallacy of taking this untrue relation of
cognition as the true relation has become the universal opinion of modern
times.
It is immediately clear from this definition of finite cognition that it is
a contradiction that sublates itself; it is the contradiction of a truth that
is supposed at the same time not to be truth, of a cognition of what is
that at the same time does not know the thing-in-itself. In the collapse of
this contradiction, its content, subjective cognition and the thing-in-itself,
collapses, that is, proves itself to be an untruth. But it is incumbent upon
cognition itself to resolve its finitude by its own forward movement and
along with it its contradiction. What we have said is a consideration which
we bring to it and remains a reflection external to it. But cognition is
itself the concept which is a purpose unto itself and, therefore, through its
realization fulfills itself, and precisely in this fulfillment sublates its subjec-
tivity and the presupposed being-in-itself. – We must examine cognition,
therefore, in its positive activity within it. Because this idea, as we have
shown, 34 is the concept’s impulse to realize itself for itself, its activity con-
sists in determining the object, and by virtue of this determining to refer
itself to itself in it as identical. The object is simply the determinable as
such, and in the idea it has this essential side of not being in and for itself
opposed to the concept. Because this cognition is still finite, not speculative
cognition, the presupposed objectivity does not as yet have for it the shape
of something which is inherently the concept simply and solely and does
not hold anything particular for itself as against the cognition. But by thus
having the status of a beyond that exists in itself, the determination of being
determinable through the concept is essential to it; for the idea is the concept
that exists for itself, is that which is absolutely infinite in itself, in which the
object is implicitly sublated, and the aim is still to sublate it explicitly. The
object, therefore, is indeed presupposed by the idea of cognition as existing
in itself, but as so essentially related to the idea that the latter, certain of
itself and of the nothingness of this opposition, arrives in the object at the
realization of its concept.
In the syllogism whereby the subjective idea now rejoins objectivity, the
first premise is the same form of immediate seizure and connection of the
concept with respect to the object as we see in the purposive connection. 35
The determining activity of the concept upon the object is an immediate
communication of itself to the object, an unresisted invasion of it. In all this
the concept remains in pure self-identity; but this immediate immanent
reflection equally has the determination of objective immediacy; that which
for the concept is its own determination, is equally a being, for it is the
first negation of the presupposition. The posited determination equally
has the status, therefore, of a presupposition which is merely found, the
apprehension of a given wherein the activity of the concept consists rather
in being negative towards itself, in holding itself back away from what is
found and passive towards it, in order that the latter be allowed to show
itself, not as determined by the subject, but as it is in itself.
In this premise, therefore, this cognition does not in any way appear as an
application of logical determinations, but as a reception and apprehension
of such determinations as already found, and its activity appears restricted
simply to the removing from the subject matter of a subjective obstacle, an
external veil. This cognition is analytic cognition.

a. Analytic cognition

The difference between analytic and synthetic cognition is sometimes said
to be that the one proceeds from the known to the unknown and the other
from the unknown to the known. On closer examination, however, it is
difficult to find any definite thought behind this difference, even less a
concept. It may be said that in general cognition begins with ignorance,
for one does not learn to know something with which one is already
acquainted. Conversely, it also begins with the known, for it is a tautological
proposition that that with which cognition begins, what it therefore actually
knows, is for that reason a known; what is as yet not known, and is expected
to be known only later, is still an unknown. In this respect it must be said
that cognition, once it has begun, always proceeds from the known to the
unknown.

The specific difference of analytic cognition is already established by the
fact that, since it is the first premise of the whole syllogism, mediation
does not as yet belong to it; analytic cognition is rather the immediate
communication of the concept, a communication that does not as yet
contain otherness and in which activity divests itself of its negativity. Yet
this immediacy of the connection is for that reason itself mediation, for it
is a negative reference of the concept to the object that annuls itself and
thereby makes itself simple and identical. This immanent reflection is only
subjective, because in its mediation the difference is present still in the
form of a presupposition existing in itself, as the object’s difference 36 within
itself. The determination that results through this connection, therefore,
is the form of simple identity, of abstract universality. Accordingly, analytic
cognition has in general this identity for its principle, and the transition
into an other, the linking of different terms is excluded from it and from
its activity.

We have said that analysis becomes synthetic when it comes to deter-
minations that are no longer posited by the problems themselves. But the
general transition from analytic to synthetic cognition lies in the necessary
transition from the form of immediacy to mediation, from abstract identity
to difference. Analysis in general restricts its activity to determinations in
so far as these are self-referential; yet by virtue of their determinateness they
also essentially refer to an other by nature. We have already remarked that
analytic cognition remains such even when it advances to relations that
are not an externally given material but are rather thought determinations,
since for it these relations are also given. But because the abstract identity
which this cognition knows to be solely its own is essentially an identity
in difference, 44 even as such it must be cognition’s own identity, and the
connection as well must become for the concept one which is posited by it
and is identical with it.

b. Synthetic cognition

Analytic cognition is the first premise of the whole syllogism – the imme-
diate reference of the concept to the object. Identity is, therefore, the
determination which analytic cognition recognizes as its own, and analytic
cognition is the apprehension of what is. Synthetic cognition aims at the
comprehension of what is, that is, at grasping the manifoldness of determi-
nations in their unity. It is, therefore, the second premise of the syllogism,
the one in which the diverse as such is connected. Its aim, therefore, is
necessity in general. – The diverse terms that are combined stand, on the
one hand, in a relation in which they are both connected yet mutually indif-
ferent and self-subsistent; but, on the other hand, they are linked together
in the concept which is their simple yet determinate unity. Now inasmuch
as in a first moment synthetic cognition passes over from abstract identity
to relation, or from being to reflection, it is not the absolute reflection
of the concept that the latter recognizes in its subject matter; the reality
that the concept gives itself is the next stage, 45 namely the said identity
in diversity as such, an identity that equally is, therefore, still inner, and
only necessity; it is not the subjective identity existing for itself, hence not
as yet the concept as such. Synthetic cognition, therefore, does also have
for its content the determinations of the concept, the object is posited in
them; but they stand only in relation to one another or in immediate unity,
and for that very reason not in the unity by which the concept exists as
subject.

This is what constitutes the finitude of this cognition. Because the iden-
tity which this real side of the idea has in it is still an inner one, the
determinations of that identity are still external to themselves; and because
the identity is not as subjectivity, the concept’s own specific presence in
the subject matter still lacks singularity; although what in the object corre-
sponds to the concept is no longer the abstract but the determinate form of
the concept and hence the concept’s particularity; the singularizing element
in the object is nevertheless still a given content. Consequently, although
this cognition transforms the objective world into concepts, what it gives
to it in accordance with conceptual determinations is only the form; as for
the object in its singularity, in its determinate determinateness, this it must
find; the cognition is not yet self-determining. It likewise finds propositions
and laws, and proves their necessity; but it proves the latter not as a necessity
inherent in a fact in and for itself, that is to say, it does not demonstrate it
from the concept; it proves it rather as the necessity inherent to a cognition
that delves into given determinations, into phenomenal differences, and
cognizes for itself the proposition as a unity and relation, or cognizes the
ground of appearance from the appearance itself.


still not realized, what we have is a relapse of the concept to the standpoint
that it assumes prior to its activity, when the actual is determined as
worthless and yet presupposed as real. This is a relapse that gives rise to the
progression to bad infinity. Its sole ground is that in the sublating of that
abstract reality the sublating itself is just as immediately forgotten, or what
is forgotten is that this reality is rather already presupposed as an actuality
which is in and for itself worthless, nothing objective. This repetition of
the presupposition of the unrealized purpose after the actual realization of
the purpose also means that the subjective attitude of the objective concept
is reproduced and perpetuated, with the result that the finitude of the good,
with respect to both content and form, appears as the abiding truth, and
its actualization always as only a singular, never universal, act. – As a matter
of fact this state has already sublated itself in the realization of the good;
what still limits the objective concept is its own view of itself, and this view
vanishes in the reflection on what its realization is in itself. By this view the
concept only stands in its own way, and all that it has to do about it is to
turn, not against an external actuality, but against itself.
That is to say, the activity in the second premise produces only a one-
sided being-for-itself, and its product therefore appears as something subjec-
tive and singular, and the first presupposition is consequently repeated in it.
But this activity is in truth just as much the positing of the implicit identity
of the objective concept and the immediate actuality. This actuality is by
presupposition determined to have only the reality of an appearance, to be
in and for itself a nullity, entirely open to determination by the objective
concept. As the external actuality is altered by the activity of the objective
concept and its determination is consequently sublated, the merely appar-
ent reality, the external determinability and worthlessness, are by that very
fact removed from it and it is thereby posited as having existence in and for
itself. In this the presupposition itself is sublated, namely the determina-
tion of the good as a merely subjective purpose restricted in content, the
necessity of first realizing it by subjective activity, and this activity itself.
In the result the mediation itself sublates itself; the result is an immedi-
acy which is not the restoration of the presupposition, but is rather the
presupposition as sublated. The idea of the concept that is determined in
and for itself is thereby posited, no longer just in the active subject but
equally as an immediate actuality; and conversely, this actuality is posited
as it is in cognition, as an objectivity that truly exists. The singularity of
the subject with which the subject was burdened by its presupposition has
vanished together with the presupposition. Thus the subject now exists
as free, universal self-identity for which the objectivity of the concept is a
given, just as immediately present to the subject as the subject immediately
knows itself to be the concept determined in and for itself. Accordingly,
in this result cognition is restored and united with the practical idea; the
previously discovered reality is at the same time determined as the realized
absolute purpose, no longer an object of investigation, a merely objective
world without the subjectivity of the concept, but as an objective world
whose inner ground and actual subsistence is rather the concept. This is
the absolute idea.

SK 37.

yasmat buddhi purusasya upabho-gam sarvam pratisadhyati
sa eva ca suksma pradhana-purusantaram visinasti

Because it is the Buddhi that accomplishes the experiences
with regard to all objects to the Purusha.
It is that again that discriminates the subtle difference
between the Pradhana and the Purusha.

In synthetic cognition, therefore, the idea achieves its purpose only to the
extent that the concept becomes for the concept according to its moments of
identity and real determinations, or of universal and particular differences –
further also as an identity which is connectedness and dependence in diversity.
But this, its subject matter, is not adequate to the concept; for in it, in this
subject matter or in its reality, the concept does not come to be the unity
of itself with itself; in necessity its identity is for it, but in this identity the
necessity is not itself the determinateness but is on the contrary a material
external to it, that is to say, is not determined by the concept and the
concept, therefore, does not recognize itself in it. Thus in general the
concept is not for itself, is not at the same time determined in and for itself
according to its unity. For this reason the idea does not as yet attain the
truth in this cognition: it does not because of the disproportion between
subject matter and subjective concept. – But the sphere of necessity is the
highest point of being and reflection; of itself, in and for itself, it passes
over into the freedom of the concept, inner identity passes over into its
manifestation which is the concept as concept. How this transition from
the sphere of necessity to the concept occurs in itself has been shown
when considering necessity, 67 and we saw it also at the beginning of this
Book as the genesis of the concept. 68 In the present context, necessity has the
position of being the reality or the subject matter of the concept, just as the
concept into which it passes is now the concept’s subject matter. But the
transition itself is the same. Here, too, it is at first only in itself, still lying
in our reflection outside cognition, that is, itself still the inner necessity
of cognition. Only the result is for cognition. The idea, in so far as the
concept is now for itself determined in and for itself, is the practical idea,
action.

B. THE IDEA OF THE GOOD

Inasmuch as the concept, which is its own subject matter, is determined in
and for itself, the subject is determined as singular. As subjective it again has
an implicit otherness for its presupposition; it is the impulse to realize itself,
the purpose that on its own wants to give itself objectivity in the objective
world and realize itself. In the theoretical idea the subjective concept, as a
universal that in and for itself lacks determination, stands opposed to the
objective world from which it derives determinate content and filling. But
in the practical idea it is as actual that it stands over against the actual;
but the certainty of itself that the subject possesses in being determined
in and for itself is a certainty of its actuality and of the non-actuality of
the world; it is the singularity of this world, and the determinateness of its
singularity, not just its otherness as abstract universality, which is a nullity
for the subject. The subject has here vindicated objectivity for itself; its
inner determinateness is the objective, for it is the universality which is just
as much absolutely determined; the previously objective world is on the
contrary only something still posited, an immediate which is determined in
a multitude of ways but which, because it is only immediately determined,
in itself eludes the unity of the concept and is of itself a nullity.
This determinateness which is in the concept, is equal to the concept,
and entails a demand for singular external actuality, is the good. It comes
on the scene with the dignity of being absolute, because it is intrinsically
the totality of the concept, the objective which is at the same time in the
form of free unity and subjectivity. This idea is superior to the idea of
cognition just considered, for it has not only the value of the universal but
also of the absolutely actual. – It is impulse, in so far as this actual is still
subjective, self-positing, without at the same time the form of immediate
presupposition; its impulse to realize itself is not, strictly speaking, to give
itself objectivity, for this it possesses within itself, but to give itself only
this empty form of immediacy. – The activity of purpose, therefore, is
not directed at itself, is not a matter of letting in a given determination
and making it its own, but of positing rather its own determination and,
by means of sublating the determinations of the external world, giving
itself reality in the form of external actuality. – The idea of the will as a
self-determining explicitly possesses content within itself. Now this content
is indeed a determinate content, and to this extent finite and restricted;
self-determination is essentially particularization, since the reflection of
the will is in itself, as negative unity as such, also singularity in the sense
that it excludes an other while presupposing it. Yet the particularity of the
content is at first infinite by virtue of the form of the concept, of which
it is the proper determinateness, and which in that content possesses its
negative self-identity, and consequently not only a particularity but its
infinite singularity. The mentioned finitude of the content in the practical
idea only means, therefore, that the idea is at first not yet realized; the
concept is for the content that which exists in and for itself; it is here the
idea in the form of objectivity existing for itself; on the one hand, the
subjective is for this reason no longer just something posited, arbitrary or
accidental, but is an absolute; but, on the other hand, this form of concrete
existence, this being-for-itself, does not as yet have the form of the being-in-
itself. Thus what from the side of the form as such appears as opposition,
appears in the form of the concept reflected into simple identity, that is,
appears in the content as its simple determinateness; the good, although
valid in and for itself, is thereby a certain particular purpose, but not one
that first receives its truth by being realized; on the contrary, it is for itself
already the true.
The syllogism of immediate realization does not itself require closer
exposition here; it is none other than the previously considered syllogism
of external purposiveness; 69 only the content constitutes the difference. In
external as in formal purposiveness it was an indeterminate finite content
in general; here, though also finite, it is as such at the same time absolutely
valid. But in regard to the conclusion, the realized purpose, a further dif-
ference enters in. In being realized the finite purpose still attains only the
status of a means; since it is not a purpose determined in and for itself
already from the beginning, as realized it also remains something that does
not exist in and for itself. If the good is again also fixed as something finite,
and is essentially such, then, notwithstanding its inner infinity, it too can-
not escape the fate of finitude – a fate that manifests itself in several forms.
The realized good is good by virtue of what it already is in the subjective
purpose, in its idea; the realization gives it an external existence, but since
this existence has only the status of an externality which is in and for itself
null, what is good in it has attained only an accidental, fragile existence,
not a realization corresponding to the idea. – Further, since this good is
restricted in content, there are several kinds of it; in concrete existence a
good is subject to destruction not only due to external contingency and
to evil, but also because of collision and conflict in the good itself. From
the side of the objective world presupposed for it (in the presupposition
of which consists the subjectivity and the finitude of the good, and which
as a distinct world runs its own course), the realization itself of the good
is exposed to obstacles, indeed, might even be made impossible. The good
thus remains an ought; it exists in and for itself, but being, as the ultimate
abstract immediacy, remains over against it also determined as a non-being.
The idea of the fulfilled good is indeed an absolute postulate, but no more
than a postulate, that is, the absolute encumbered with the determinate-
ness of subjectivity. There still are two worlds in opposition, one a realm of
subjectivity in the pure spaces of transparent thought, the other a realm of
objectivity in the element of an externally manifold actuality, an impervious
realm of darkness. The complete development of this unresolved contra-
diction, between that absolute purpose and the restriction of this reality that
stands opposed to it, has been examined in detail in the Phenomenology
of Spirit (pp. 323ff.). 70 – Inasmuch as the idea has within it the moment
of complete determinateness, the other concept to which the concept in
it relates possesses in its subjectivity at the same time the moment of an
object; consequently the idea enters here into the shape of self-consciousness,
and in this one respect coincides with its exposition.
But what the practical idea still lacks is the moment of real consciousness
itself, namely that the moment of actuality in the concept would have
attained for itself the determination of external being. – This lack can also
be regarded in this way, namely that the practical idea still lacks the moment
of the theoretical idea. That is to say, in the latter there stands on the side
of the subjective concept – the concept that is in process of being intuited
in itself by the concept – only the determination of universality; cognition
only knows itself as apprehension, as the identity of the concept with
itself which, for itself, is indeterminate; the filling, that is, the objectivity
determined in and for itself, is for this identity a given; what truly exists is for
it the actuality present there independently of any subjective positing. For
the practical idea, on the contrary, this actuality constantly confronting it
as an insuperable restriction is in and for itself a nullity that ought to receive
its true determination and intrinsic value only through the purposes of the
good. It is the will, therefore, that alone stands in the way of attaining its
goal, because it separates itself from cognition and because for it external
actuality does not receive the form of a true existence. The idea of the good
can therefore find its completion only in the idea of the true.
But it makes this transition through itself. In the syllogism of action, one
premise is the immediate reference of the good purpose to the actuality which
it appropriates and which, in the second premise, it directs as external
means against the external actuality. The good is for the subjective concept
the objective; actuality confronts it in existence as an insuperable restriction
only in so far as it still has the determination of immediate existence, not
of something objective in the sense that it is being in and for itself; it is
rather either the evil or the indifferent, the merely determinable, whose
worth does not lie within it. But this abstract being that confronts the
good in the second premise has already been sublated by the practical idea
itself; the first premise of this idea’s action is the immediate objectivity of
the concept, according to which purpose is communicated to actuality
without any resistance and is in the simple connection of identity with it.
To this extent, therefore, what remains is to bring together the thoughts of
the two premises of the practical idea. All that is added to what is already
accomplished in the first premise by the objective concept is that in the
second it is posited by way of mediation, hence for it. Just as in purposive
connection in general, where the realized purpose is again only a means
but the means is conversely also the realized purpose, so too now in the
syllogism of the good the second premise is already immediately present in
the first in itself, except that this immediacy is not sufficient and the second
premise is for the first already postulated – the realization of the good in the
face of another actuality confronting it is the mediation which is essentially
necessary for the immediate connection and consummation of the good.
For the first premise is only the first negation or the otherness of the concept,
an objectivity that would be a state of immersion of the concept into
externality; the second premise is the sublation of this otherness, whereby
the immediate realization of the purpose first becomes the actuality of
the good as concept existing for itself, for in that actuality the concept is
posited as identical with itself, not with an other, and in this way alone as
free concept. If it is now claimed that the purpose of the good is thereby
still not realized, what we have is a relapse of the concept to the standpoint
that it assumes prior to its activity, when the actual is determined as
worthless and yet presupposed as real. This is a relapse that gives rise to the
progression to bad infinity. Its sole ground is that in the sublating of that
abstract reality the sublating itself is just as immediately forgotten, or what
is forgotten is that this reality is rather already presupposed as an actuality
which is in and for itself worthless, nothing objective. This repetition of
the presupposition of the unrealized purpose after the actual realization of
the purpose also means that the subjective attitude of the objective concept
is reproduced and perpetuated, with the result that the finitude of the good,
with respect to both content and form, appears as the abiding truth, and
its actualization always as only a singular, never universal, act. – As a matter
of fact this state has already sublated itself in the realization of the good;
what still limits the objective concept is its own view of itself, and this view
vanishes in the reflection on what its realization is in itself. By this view the
concept only stands in its own way, and all that it has to do about it is to
turn, not against an external actuality, but against itself.
That is to say, the activity in the second premise produces only a one-
sided being-for-itself, and its product therefore appears as something subjec-
tive and singular, and the first presupposition is consequently repeated in it.
But this activity is in truth just as much the positing of the implicit identity
of the objective concept and the immediate actuality. This actuality is by
presupposition determined to have only the reality of an appearance, to be
in and for itself a nullity, entirely open to determination by the objective
concept. As the external actuality is altered by the activity of the objective
concept and its determination is consequently sublated, the merely appar-
ent reality, the external determinability and worthlessness, are by that very
fact removed from it and it is thereby posited as having existence in and for
itself. In this the presupposition itself is sublated, namely the determina-
tion of the good as a merely subjective purpose restricted in content, the
necessity of first realizing it by subjective activity, and this activity itself.
In the result the mediation itself sublates itself; the result is an immedi-
acy which is not the restoration of the presupposition, but is rather the
presupposition as sublated. The idea of the concept that is determined in
and for itself is thereby posited, no longer just in the active subject but
equally as an immediate actuality; and conversely, this actuality is posited
as it is in cognition, as an objectivity that truly exists. The singularity of
the subject with which the subject was burdened by its presupposition has
vanished together with the presupposition. Thus the subject now exists
as free, universal self-identity for which the objectivity of the concept is a
given, just as immediately present to the subject

YS III.21

    kaya-rupa-samyamat tad-grahya-sakti-stambhe caksu-prakasa-asamprayoge ‘ntardhanam

YS III.22

    etena sabda-adi-antardhanam uktam

The absolute idea has shown itself to be the identity of the theoretical and
the practical idea, each of which, of itself still one-sided, possesses the idea
only as a sought-for beyond and unattained goal; each is therefore a synthesis
of striving, each possessing as well as not possessing the idea within it, passing
over from one thought to the other without bringing the two together but
remaining fixed in the contradiction of the two. The absolute idea, as the
rational concept that in its reality only rejoins itself, is by virtue of this
immediacy of its objective identity, on the one hand, a turning back to life;
on the other hand, it has equally sublated this form of its immediacy and
harbors the most extreme opposition within. The concept is not only soul,
but free subjective concept that exists for itself and therefore has personality –
the practical objective concept that is determined in and for itself and
is as person impenetrable, atomic subjectivity – but which is not, just the
same, exclusive singularity; it is rather explicitly universality and cognition,
and in its other has its own objectivity for its subject matter. All the rest
is error, confusion, opinion, striving, arbitrariness, and transitoriness; the
absolute idea alone is being, imperishable life, self-knowing truth, and is all
truth.

It is the sole subject matter and content of philosophy. Since it contains
all determinateness within it, and its essence consists in returning through
its self-determination and particularization back to itself, it has various
shapes, and the business of philosophy is to recognize it in these. Nature
and spirit are in general different modes of exhibiting its existence, art
and religion its different modes of apprehending itself and giving itself
appropriate existence. Philosophy has the same content and the same
purpose as art and religion, but it is the highest mode of apprehending the
absolute idea, because its mode, that of the concept, is the highest. Hence
it seizes those shapes of real and ideal finitude, as well of infinity and
holiness, and comprehends them and itself. The derivation and cognition
of these particular modes are now the further business of the particular
philosophical sciences. Also the logicality of the absolute idea can be called
a mode of it; but mode signifies a particular kind, a determinateness of
form, whereas the logicality of the idea is the universal mode in which
all particular modes are sublated and enveloped. The logical idea is the
idea itself in its pure essence, the idea which is enclosed in simple identity
within its concept and in reflective shining has as yet to step into a form-
determinateness. The Logic thus exhibits the self-movement of the absolute
idea only as the original word, a word which is an utterance, but one that
in being externally uttered has immediately vanished again. The idea is,
therefore, only in this self-determination of apprehending itself; it is in pure
thought, where difference is not yet otherness, but is and remains perfectly
transparent to itself. – The logical idea thus has itself, as the infinite form,
for its content – form that constitutes the opposite of content inasmuch as
the latter is the form determination that has withdrawn into itself and has
been so sublated in identity that this concrete identity stands over against
the identity developed as form; the content has the shape of an other and of
something given as against the form that as such stands simply in reference,
and whose determinateness is posited at the same time as reflective shine. –
More exactly, the absolute idea itself has only this for its content, namely
that the form determination is its own completed totality, the pure content.
Now the determinateness of the idea and the entire course traversed by this
determinateness has constituted the subject matter of the science of logic,
and out of this course the absolute idea has come forth for itself; thus to
be for itself, however, has shown itself to amount to this, namely that
determinateness does not have the shape of a content, but that it is simply
as form, and that accordingly the idea is the absolutely universal idea. What
is left to be considered here, therefore, is thus not a content as such, but
the universal character of its form – that is, method.
Method may appear at first to be just the manner in which cognition
proceeds, and this is in fact its nature. But as method this manner of
proceeding is not only a modality of being determined in and for itself; it is a
modality of cognition, and as such is posited as determined by the concept
and as form, since form is the soul of all objectivity and all otherwise
determined content has its truth in form alone. If the content is again
assumed as given to the method and of a nature of its own, then method,
so understood, is just like the logical realm in general a merely external
form. But against this assumption appeal can be made, not only to the
fundamental concept of what constitutes logic, but to the entire logical
course in which all the shapes of a given content and of objects came up for
consideration. This course has shown the transitoriness and the untruth of
all such shapes; also that no given object is capable of being the foundation
to which the absolute form would relate as only an external and accidental
determination; that, on the contrary, it is the absolute form that has proved
itself to be the absolute foundation and the ultimate truth. For this course
the method has resulted as the absolutely self-knowing concept, as the concept
that has the absolute, both as subjective and objective, as its subject matter,
and consequently as the pure correspondence of the concept and its reality,
a concrete existence that is the concept itself.
Accordingly, what is to be considered as method here is only the move-
ment of the concept itself. We already know the nature of this movement,
but it now has, first, the added significance that the concept is all, and
that its movement is the universal absolute activity, the self-determining
and self-realizing movement. The method is therefore to be acknowledged
as the universal, internal and external mode, free of restrictions, and as
the absolutely infinite force to which no object that may present itself as
something external, removed from reason and independent of it, could
offer resistance, or be of a particular nature opposite to it, and could not be
penetrated by it. It is therefore soul and substance, and nothing is conceived
and known in its truth unless completely subjugated to the method; it is the
method proper to each and every fact because its activity is the concept.
This is also the truer meaning of its universality; according to the universal-
ity of reflection, it is taken only as the method for all things; but according
to the universality of the idea, it is both the manner of cognition, of the
concept subjectively aware of itself, and the objective manner, or rather the
substantiality of things – that is, of concepts as they first appear as others
to representation and reflection. It is therefore not only the highest force of
reason, or rather its sole and absolute force, but also reason’s highest and
sole impulse to find and recognize itself through itself in all things. – Second,
here we also have the distinction of the method from the concept as such, the
particularization of the method. As the concept was considered for itself,
it appeared in its immediacy; the reflection, or the concept considering it,
fell on the side of our knowledge. The method is this knowledge itself, for
which the concept is not only as subject matter but is as its own subjec-
tive act, the instrument and the means of cognitive activity, distinct from
this activity and yet the activity’s own essentiality. In cognition as enquiry,
the method likewise occupies the position of an instrument, as a means
that stands on the side of the subject, connecting it with the object. The
subject in this syllogism is one extreme, the object is the other, and in
conclusion the subject unites through its method with the object without
however uniting with itself there. The extremes remain diverse, because
subject, method, and object are not posited as the one identical concept; the
syllogism is therefore always the formal syllogism; the premise in which
the subject posits the form on its side as its method is an immediate deter-
mination and contains therefore the determinations of the form – as we
have seen, of definition, division, and so forth 71 – as matters of fact found
ready-made in the subject. In true cognition, on the contrary, method is
not only an aggregate of certain determinations, but the determinateness
in-and-for-itself of the concept, and the concept is the middle term only
because it equally has the significance of the objective; in the conclusion,
therefore, the objective does not attain only an external determinateness by
virtue of the method, but is posited rather in its identity with the subjective
concept.
